\subsection{Directed local PL limits retain the actual curve maximum argument}

\begin{lemma}[Actual curve-piece traces of directed local PL limits]
\label{lem:fc-directed-pl-curve-piece-trace}
\lean{EnvelopeProperty.Berkovich.RelativeSpectrum.FCStandardAffinoid.isLineSubharmonicOn_curveRelativePieceNorm_of_isDecrConvexPLNetLimitOn,EnvelopeProperty.Berkovich.RelativeSpectrum.FCStandardAffinoid.isLineSubharmonicOn_curveRelativePieceNorm_of_isLocallyNetPSHOn}
\leanok
A decreasing directed infimum of original regular-unit convex PL functions
has subharmonic finite multiplicity-weighted norm on every actual selected
finite-flat curve piece. The assertion is local on the original curve.
\end{lemma}
\begin{proof}
Each original convex PL term has the proved source line-comparison norm.
On each actual finite fibre, addition and multiplication by the original
positive multiplicities commute with directed decreasing infima; no term
reaches plus infinity. Thus the unchanged whole finite norm is the directed
infimum of the original comparison norms. Source line subharmonicity is closed
under such directed infima, using the finite actual boundary to choose a
single common tail index for any strict comparison bound. For locally given
presentations, refine around the finite central fibre into disjoint actual
idempotent pieces, each contained in its own presentation neighborhood. Apply
the global-presentation argument on each piece, sum their actual traces, and
use the exact restriction and partition identities to recover the original
whole piece norm. Source line locality gives the assertion.
\end{proof}

\begin{theorem}[Actual local maximum of directed locally PL limits]
\label{thm:fc-directed-pl-curve-local-maximum}
\lean{EnvelopeProperty.Berkovich.RelativeSpectrum.IsLocallyNetPSHOn.eventually_eq_of_isLocalMax_finiteFree}
\leanok
\uses{lem:fc-directed-pl-curve-piece-trace}
On a genuine finite-free algebraic curve over an algebraically closed field
with a nontrivial valuation, a local maximum of a function locally given by
decreasing directed original convex PL nets is locally constant.
\end{theorem}
\begin{proof}
If the maximum is minus infinity, the maximum inequality gives the assertion.
Otherwise write it as a finite real number $R$. Construct an actual single-point
central idempotent piece inside a presentation neighborhood and the actual
local sublevel neighborhood $u\le R$. Its unchanged finite norm is subharmonic
by the preceding lemma, bounded above by the full original piece weight times
$R$, and equals this bound at the central base point. The source line local
maximum principle makes that norm constant nearby. Every original fibre
multiplicity is positive and the full weight remains constant, so equality
of the finite weighted sum forces every actual summand to equal $R$. The
actual selected piece and base interior are neighborhoods, giving equality
on an original neighborhood of the original maximizing point.
\end{proof}

\begin{lemma}[Every original valuation retains the directed local maximum]
\label{lem:fc-directed-pl-all-valuations-local-maximum}
\lean{EnvelopeProperty.Berkovich.RelativeSpectrum.IsLocallyNetPSHOn.eventually_eq_of_isLocalMax_finiteFree_allValuations}
\leanok
\uses{thm:fc-directed-pl-curve-local-maximum}
The preceding finite-free local maximum statement holds for every original
field valuation, including the trivial valuation.
\end{lemma}
\begin{proof}
Construct the actual nontrivially valued complete algebraically closed scalar
extension of the prescribed original valuation. Scalar extension of the
original finite-free algebra stays finite-free over the literal polynomial
scalar extension, identified with the new polynomial ring. Pull each original
regular-unit directed presentation through actual scalar restriction. Every
native point above the original local maximizing point is then a local
maximizer of the pulled-back function, so the preceding theorem gives a
plateau near each such lift. The genuine whole native scalar restriction is
closed and surjective. Closedness gives an original neighborhood whose whole
inverse image lies in the union of these plateau neighborhoods. Surjectivity
then lifts each point of this original neighborhood and proves the unchanged
original equality there.
\end{proof}

\begin{theorem}[Actual integral scheme-curve directed maximum principle]
\label{thm:fc-integral-curve-directed-maximum}
\lean{EnvelopeProperty.Berkovich.RelativeSchemeValuedPoint.IsLocallyNetPSHOn.eq_of_isMaxOn_univ_integral_curve}
\leanok
\uses{lem:fc-directed-pl-all-valuations-local-maximum}
On an original integral scheme of finite type and global dimension at most
one over a complete algebraically closed ultrametric normed field, a function
locally given by decreasing directed original regular-unit convex PL nets is
constant if it attains its global maximum.
\end{theorem}
\begin{proof}
On each nonempty affine chart the original coordinate ring is a finite-type
domain of dimension at most one. If it is a field, algebraic closedness makes
the entire native point space a singleton. Otherwise actual finite curve
normalization supplies a finite-free polynomial parameter. The preceding
all-valuations theorem proves local maximum plateaus on this unchanged affine
chart, and the genuine native affine open embedding transports them to the
original scheme.

The actual whole native integral curve is preconnected by its proved original
geometric connectedness. Directed local PL limits are upper semicontinuous
and never attain plus infinity, because these assertions hold for each
original convex PL term and are retained by pointwise infima and locality.
The set of points attaining the original global maximum is therefore closed.
The local plateau statement makes it open. It contains the original maximizing
point, so preconnectedness makes it the whole original native curve.
\end{proof}

\begin{theorem}[Proper directed native maximum over the original field]
\label{thm:fc-proper-directed-native-maximum}
\lean{EnvelopeProperty.Berkovich.RelativeSchemeValuedPoint.IsLocallyNetPSHOn.eq_of_isMaxOn_univ_geometricallyConnected_proper}
\leanok
\uses{thm:fc-integral-curve-directed-maximum,lem:fc-native-complete-algebraic-closure-surjective,thm:fc-connected-proper-native-maximum}
On an actual proper geometrically connected scheme over any complete
ultrametric normed field, a function locally given by decreasing directed
original regular-unit convex PL nets is constant if it attains a global
maximum. No algebraic-closedness or nontrivial valuation hypothesis is added.
\end{theorem}
\begin{proof}
Over an algebraically closed base, first apply the integral curve directed
maximum principle to each actual closed integral curve of the genuine
closed-intermediate finite-type curve chain. This propagates a maximum between
all original closed supports. For an arbitrary support, take its actual
integral reduced singleton closure and an actual original rational section.
The integral proper coefficient theorem makes this closure geometrically
connected; rationalize its original native point over the actual completed
algebraic closure of its completed residue field. That rational lift and the
scalar-changed original section have closed supports on an actually connected
proper scalar-changed source. Closed-support propagation equates their values.
First do this at the maximizing support, then at every other support, comparing
the resulting original closed rational anchors by the same original curve
chains. Thus the function is constant on the whole original native space.

For an arbitrary complete base take its actual completed algebraic closure.
The literal proper scheme pullback is connected by the original geometric
connectedness. Every original point lifts through the proved whole native
scalar restriction. Directed regular-unit presentations pull back term by
term, and the lifted original maximizer is an actual global maximizer. The
algebraically closed conclusion equates the two lifted values. Their exact
original native projections recover the original values, proving the result.
\end{proof}

\begin{remark}[Directed comparison and the original sequence class]
\label{rmk:fc-directed-maximum-original-class}
The directed class is kept distinct from the original sequence-based locally
psh definition. The proof uses its actual multiplicity-weighted curve norm
comparisons and never extracts a cofinal sequence. These stronger comparison
statements allow the original global psh directed approximants to prove fibre
constancy directly, while the source's separate unrestricted sequential
regularization statement remains a separate formalization obligation.
\end{remark}
